\section{Relativity}

Relativity is not an assumption of the model, it is a large-scale consequence of the mesh having no preferred
direction. The twenty-four directions of a dock are so symmetric that, at large scale, the speed limit and the way space
and time trade off, Lorentz invariance, emerge from that even-handedness.

There is a tiny residual lumpiness, because the mesh is discrete, but it is governed by the high symmetry of the cell
and enters only at sixth order in the energy over the Planck scale. That is about one part in ten to the hundred and
nine at the highest gamma-ray-burst energies, unobservably small, and far below the first-order violation a flat lattice
would give, which the gamma-ray-burst data already rule out. So the substrate passes the sharpest test of Lorentz
invariance comfortably where a flat mesh fails.

The deeper reason is the one the physical-space section gave. The cusp is an aperiodic flat slice with no repeating
grain, so its roughness is random and self-averages away the farther out you look, restoring isotropy with scale, while
a periodic lattice keeps its facets at every scale and stays anisotropic. Isotropy is not put in, it is what a curved
aperiodic substrate flows to and a flat mesh cannot, which is exactly why a relativistic world needs the curved
substrate and not the flat one.
